%% ============================================================
%%  Aula 2 — Tomando boas decisões dado um modelo do mundo
%%  Curso: Aprendizagem por Reforço — UFU
%%  Prof. Marcelo Keese Albertini
%%  Adaptado e expandido de CS234 (Stanford), Emma Brunskill
%%  COMPILAR COM XeLaTeX (duas vezes)
%% ============================================================
\documentclass[aspectratio=169, 11pt]{beamer}

\usetheme{AprendizadoRL}      % ANTES do babel
\usepackage{babel}
\babelprovide[import, main]{brazilian}
\usepackage{booktabs}

\title[Aula 2 — Decisões com modelo]{Aula 2: Tomando sequências de boas decisões\\ dado um modelo do mundo}
\subtitle[Aprendizagem por Reforço]{Processos de decisão de Markov \textbullet\ Avaliação \textbullet\ Iteração de política e de valor}
\author[Marcelo K. Albertini]{Prof. Marcelo Keese Albertini \\
  \footnotesize Universidade Federal de Uberlândia \\
  \footnotesize adaptado de Emma Brunskill, CS234 (Stanford)}
\institute{Aprendizagem por Reforço \textbullet\ UFU}
\date{2026}

%% ------------------------------------------------------------
%% Macros locais desta aula
%% ------------------------------------------------------------
\newcommand{\Bop}{\mathcal{B}}                 % operador de backup de Bellman
\newcommand{\Bpol}{\mathcal{B}^{\pol}}         % operador para política fixa
\newcommand{\Qstar}{Q^{*}}
\newcommand{\Rmax}{R_{\max}}

%% Mars Rover determinístico: 7 estados, ações a1 (esquerda) e a2 (direita)
\newcommand{\marsroverdet}{%
\begin{tikzpicture}[node distance=6.5mm]
  \foreach \i [count=\c] in {1,...,7}{%
    \ifnum\c=1
      \node[estado] (s\i) {$s_{\i}$};
    \else
      \pgfmathtruncatemacro{\prev}{\i-1}
      \node[estado, right=of s\prev] (s\i) {$s_{\i}$};
    \fi}
  \node[estado destacado] at (s1) {$s_1$};
  \node[estado destacado] at (s7) {$s_7$};
  \foreach \i in {1,...,6}{%
    \pgfmathtruncatemacro{\nxt}{\i+1}
    \draw[trans destac] (s\i) to[bend left=42] (s\nxt);}
  \draw[trans destac] (s7) to[out=35, in=-35, looseness=6] (s7);
  \foreach \i in {2,...,7}{%
    \pgfmathtruncatemacro{\prv}{\i-1}
    \draw[trans] (s\i) to[bend left=42] (s\prv);}
  \draw[trans] (s1) to[out=145, in=215, looseness=6] (s1);
  \node[font=\scriptsize\bfseries, color=rlLaranja, above=3.6mm of s4] {$a_2$: direita};
  \node[font=\scriptsize\bfseries, color=rlCinza,   below=3.6mm of s4] {$a_1$: esquerda};
  \node[recomp, above=3.6mm of s1] {$+1$};
  \node[recomp, above=3.6mm of s7] {$+10$};
\end{tikzpicture}}

%% Mars Rover estocástico (passeio aleatório)
\newcommand{\marsroverestoc}{%
\begin{tikzpicture}[node distance=9mm]
  \foreach \i [count=\c] in {1,...,7}{%
    \ifnum\c=1
      \node[estado] (s\i) {$s_{\i}$};
    \else
      \pgfmathtruncatemacro{\prev}{\i-1}
      \node[estado, right=of s\prev] (s\i) {$s_{\i}$};
    \fi}
  \node[estado destacado] at (s1) {$s_1$};
  \node[estado destacado] at (s7) {$s_7$};
  \foreach \i in {1,...,6}{%
    \pgfmathtruncatemacro{\nxt}{\i+1}
    \draw[trans] (s\i) to[bend left=40] node[probl] {0.4} (s\nxt);
    \draw[trans] (s\nxt) to[bend left=40] node[probl] {0.4} (s\i);}
  \draw[trans] (s1) to[out=-125, in=-55, looseness=7] node[probl] {0.6} (s1);
  \draw[trans] (s7) to[out=-125, in=-55, looseness=7] node[probl] {0.6} (s7);
  \foreach \i in {2,...,6}{%
    \draw[trans] (s\i) to[out=-125, in=-55, looseness=7] node[probl] {0.2} (s\i);}
  \node[recomp, above=6mm of s1] {$+1$};
  \node[recomp, above=6mm of s7] {$+10$};
\end{tikzpicture}}

%% Diagrama de backup com maximização
\newcommand{\backupotimo}{%
\begin{tikzpicture}[grow=down, level distance=10mm,
    level 1/.style={sibling distance=24mm},
    level 2/.style={sibling distance=11mm}]
  \node[estado destacado] {$s$}
    child {node[acaon] {}
      child {node[estado] {} edge from parent[trans]}
      child {node[estado] {} edge from parent[trans]}
      edge from parent[trans destac] node[left, font=\scriptsize, color=rlLaranja] {$a_1$}}
    child {node[acaon] {}
      child {node[estado] {} edge from parent[trans]}
      child {node[estado] {} edge from parent[trans]}
      edge from parent[trans] node[right, font=\scriptsize, color=rlCinza] {$a_2$}};
\end{tikzpicture}}

%% Ciclo de iteração de política generalizada
\newcommand{\ciclogpi}{%
\begin{tikzpicture}
  \node[draw=rlAzulClaro, very thick, rounded corners=2pt, fill=rlAzulClaro!8,
        text width=2.9cm, align=center, inner sep=4pt, font=\small] (av)
        {\textbf{Avaliação}\\[1pt] {\scriptsize $V \rightarrow \Vpi$}};
  \node[draw=rlLaranja, very thick, rounded corners=2pt, fill=rlLaranja!8,
        text width=2.9cm, align=center, inner sep=4pt, font=\small,
        right=2.4cm of av] (me)
        {\textbf{Melhoria}\\[1pt] {\scriptsize $\pol \rightarrow \text{gulosa}(V)$}};
  \draw[trans destac] (av) to[bend left=30] (me);
  \draw[trans destac] (me) to[bend left=30] (av);
\end{tikzpicture}}

\begin{document}

%% ============================================================
\begin{frame}[plain, noframenumbering]\titlepage\end{frame}

%% ------------------------------------------------------------
\begin{frame}{Onde estamos no curso}
  \begin{columns}[T]
    \begin{column}{0.32\textwidth}
      {\color{rlCinza}\small\bfseries AULA ANTERIOR}\\[0.3em]
      \begin{itemize}\setlength{\itemsep}{0.3ex}
        \item Introdução à aprendizagem por reforço
        \item Componentes de um agente: \alert{modelo}, \alert{valor}, \alert{política}
        \item Processos de Markov
      \end{itemize}
    \end{column}
    \begin{column}{0.34\textwidth}
      {\color{rlLaranja}\small\bfseries HOJE}\\[0.3em]
      \begin{itemize}\setlength{\itemsep}{0.3ex}
        \item Decidir bem \emph{quando o modelo do mundo é conhecido}
        \item Avaliação e controle em MDPs
        \item Iteração de política e de valor
      \end{itemize}
    \end{column}
    \begin{column}{0.32\textwidth}
      {\color{rlCinza}\small\bfseries PRÓXIMA AULA}\\[0.3em]
      \begin{itemize}\setlength{\itemsep}{0.3ex}
        \item Avaliação de política \emph{sem} modelo do mundo
        \item Monte Carlo e diferenças temporais
      \end{itemize}
    \end{column}
  \end{columns}
  \vfill
  \begin{ideiabox}
    Hoje é o único momento do curso em que \emph{tudo} é conhecido. O problema deixa de
    ser de aprendizagem e passa a ser de \alert{planejamento} --- e é por isso que ele
    fornece o alicerce teórico de todo o resto.
  \end{ideiabox}
\end{frame}

%% ------------------------------------------------------------
\begin{frame}{Roteiro}
  \begin{columns}[T]
    \begin{column}{0.55\textwidth}
      \tableofcontents[hideallsubsections]
    \end{column}
    \begin{column}{0.43\textwidth}
      \begin{defbox}{os dois problemas}
        Dado o modelo $(\gS,\gA,P,\rec,\gamma)$:
        \begin{enumerate}\setlength{\itemsep}{0.2ex}
          \item \textbf{Predição} --- qual o valor de uma política dada?
          \item \textbf{Controle} --- qual a melhor política possível?
        \end{enumerate}
      \end{defbox}
    \end{column}
  \end{columns}
\end{frame}

%% ------------------------------------------------------------
\begin{frame}{Duas perguntas para começar}
  \begin{quizbox}{}
    Em um processo de decisão de Markov, um fator de desconto $\gamma$ \alert{grande}
    significa que recompensas de \emph{curto} prazo são muito mais influentes do que as
    de longo prazo.
  \end{quizbox}
  \pause
  \begin{respbox}
    \textbf{Falso.} Um $\gamma$ grande dá \emph{mais} peso a recompensas distantes.
    No extremo, $\gamma=0$ produz um agente míope; $\gamma \to 1$, um agente que trata
    todas as épocas quase igualmente.
  \end{respbox}
  \begin{ideiabox}
    \textbf{Fio condutor da aula.} É possível construir algoritmos tais que, a cada
    iteração de computação, a política melhore \alert{monotonicamente}? Todos os
    algoritmos têm essa propriedade?
  \end{ideiabox}
\end{frame}

%% ============================================================
\section{Do processo de Markov ao MDP}
%% ============================================================

\begin{frame}{Recapitulação: a hierarquia de modelos}
  \begin{columns}[T]
    \begin{column}{0.5\textwidth}
      \begin{defbox}{propriedade de Markov}
        O futuro é independente do passado dado o presente:
        \[
          \Prob(s_{t+1} \mid s_t, a_t, s_{t-1}, \dots) = \Prob(s_{t+1} \mid s_t, a_t).
        \]
        O estado é uma \emph{estatística suficiente} da história.
      \end{defbox}
      \begin{alertabox}
        Markov não é propriedade do mundo, e sim da \alert{representação de estado}
        escolhida.
      \end{alertabox}
    \end{column}
    \begin{column}{0.48\textwidth}
      \begin{itemize}\setlength{\itemsep}{0.7ex}
        \item \textbf{Processo de Markov} $(\gS, P)$ \\
              {\small só aleatoriedade}
        \item \textbf{Processo de recompensa de Markov} $(\gS, P, \rec, \gamma)$ \\
              {\small \alert{+ recompensa}: já dá para \emph{avaliar}}
        \item \textbf{Processo de decisão de Markov} $(\gS, \gA, P, \rec, \gamma)$ \\
              {\small \alert{+ ação}: agora há o que \emph{decidir}}
      \end{itemize}
      \vspace{0.5em}
      \begin{ideiabox}
        Cada nível acrescenta um ingrediente --- e reaproveita os algoritmos do anterior.
      \end{ideiabox}
    \end{column}
  \end{columns}
\end{frame}

%% ------------------------------------------------------------
\begin{frame}{Horizonte, retorno e função valor}
  \begin{columns}[T]
    \begin{column}{0.52\textwidth}
      \begin{defbox}[fontupper=\small]{horizonte $H$}
        Número de passos por episódio. Pode ser infinito; caso contrário, o processo é
        de \emph{horizonte finito}.
      \end{defbox}
      \begin{defbox}[fontupper=\small]{retorno $G_t$}
        \vspace{-1.1em}
        \[
          G_t = r_t + \gamma r_{t+1} + \gamma^2 r_{t+2} + \cdots + \gamma^{H-1} r_{t+H-1}
        \]
        \vspace{-1.6em}
      \end{defbox}
      \begin{defbox}[fontupper=\small]{função valor $V(s)$}
        \vspace{-1.1em}
        \[ V(s) = \E[\, G_t \mid s_t = s \,] \]
        \vspace{-1.6em}
      \end{defbox}
    \end{column}
    \begin{column}{0.46\textwidth}
      \begin{ideiabox}
        $G_t$ é variável \alert{aleatória} --- muda a cada trajetória. $V(s)$ é um
        \alert{número}: a esperança já eliminou o acaso. Boa parte de RL consiste em
        estimar $V$ a partir de amostras de $G_t$.
      \end{ideiabox}
      \begin{alertabox}
        Duas políticas com o mesmo $V(s)$ podem ter distribuições de $G_t$ muito
        diferentes: a média ignora o risco.
      \end{alertabox}
    \end{column}
  \end{columns}
\end{frame}

%% ------------------------------------------------------------
\begin{frame}{Por que descontar o futuro?}
  \begin{columns}[T]
    \begin{column}{0.55\textwidth}
      \begin{enumerate}\setlength{\itemsep}{0.7ex}
        \item \textbf{Matemática:} com $\abs{r} \le \Rmax$ e $\gamma<1$, a série
              converge e $\abs{G_t} \le \Rmax/(1-\gamma)$.
        \item \textbf{Modelagem:} $\gamma$ é a probabilidade de o episódio
              \emph{continuar}.
        \item \textbf{Economia:} recompensa hoje vale mais do que amanhã.
        \item \textbf{Algorítmica:} é $\gamma<1$ que torna o operador de Bellman uma
              \alert{contração} --- fonte de todas as garantias de hoje.
      \end{enumerate}
    \end{column}
    \begin{column}{0.43\textwidth}
      \begin{ideiabox}
        Horizonte efetivo $\approx \dfrac{1}{1-\gamma}$ passos.
        \smallskip

        $\gamma = 0{,}9 \Rightarrow \sim\!10$ \qquad
        $\gamma = 0{,}99 \Rightarrow \sim\!100$
      \end{ideiabox}
      \begin{alertabox}
        $\gamma = 1$ só é seguro com horizonte finito, ou se um estado terminal for
        atingido com probabilidade $1$.
      \end{alertabox}
    \end{column}
  \end{columns}
\end{frame}

%% ------------------------------------------------------------
\begin{frame}{Equação de Bellman para um MRP}{A propriedade de Markov permite recursão}
  \[
    V(s) \;=\; \termo{\rec(s)}{recompensa imediata}
    \;+\; \termo{\gamma \sum_{s' \in \gS} \tran{s'}{s}\, V(s')}{soma descontada do futuro}
  \]
  \vspace{0.4em}
  \begin{ideiabox}
    Todo o valor de uma trajetória infinita cabe em uma equação com \emph{um} passo de
    antecipação. O valor de hoje se decompõe no que ganho agora mais o valor de onde
    posso parar.
  \end{ideiabox}
  \begin{alertabox}
    Isso é um \alert{sistema} de $N = \abs{\gS}$ equações acopladas --- uma por estado ---
    e não uma fórmula fechada. Resolvê-lo é o assunto dos próximos slides.
  \end{alertabox}
\end{frame}

%% ------------------------------------------------------------
\begin{frame}{Forma matricial e solução analítica}
  Empilhando os $N$ estados: $V, \rec \in \R^N$ e $P \in \R^{N \times N}$.
  \vspace{-0.2em}
  \[
    \begin{bmatrix} V(s_1) \\ \vdots \\ V(s_N) \end{bmatrix}
    =
    \begin{bmatrix} \rec(s_1) \\ \vdots \\ \rec(s_N) \end{bmatrix}
    + \gamma
    \begin{bmatrix}
      \tran{s_1}{s_1} & \cdots & \tran{s_N}{s_1} \\
      \vdots & \ddots & \vdots \\
      \tran{s_1}{s_N} & \cdots & \tran{s_N}{s_N}
    \end{bmatrix}
    \begin{bmatrix} V(s_1) \\ \vdots \\ V(s_N) \end{bmatrix}
  \]
  \vspace{-0.6em}
  \begin{columns}[c]
    \begin{column}{0.44\textwidth}
      \[
        \begin{aligned}
          V &= \rec + \gamma P V \\
          (I - \gamma P)\, V &= \rec \\[0.3em]
          \mdestaque{rlLaranja}{V = (I - \gamma P)^{-1} \rec}
        \end{aligned}
      \]
    \end{column}
    \begin{column}{0.54\textwidth}
      \begin{teobox}{a inversa sempre existe se $\gamma<1$}
        $\norminf{\gamma P} = \gamma < 1$, logo a série de Neumann converge:
        $(I-\gamma P)^{-1} = \sum_{k \ge 0} (\gamma P)^k$.
        O termo $k$ é a contribuição das recompensas $k$ passos à frente.
      \end{teobox}
      {\small Custo: $\Theta(N^3)$ --- viável só para $N$ moderado.}
    \end{column}
  \end{columns}
\end{frame}

%% ------------------------------------------------------------
\begin{frame}{Algoritmo iterativo para o valor de um MRP}
  \begin{algbox}{avaliação iterativa (programação dinâmica)}
    \begin{itemize}\setlength{\itemsep}{0.2ex}
      \item Inicialize $V_0(s) = 0$ para todo $s$
      \item Para $k = 1, 2, \dots$ até convergir, e para todo $s \in \gS$:
        \[ V_k(s) \leftarrow \rec(s) + \gamma \sum_{s' \in \gS} \tran{s'}{s}\, V_{k-1}(s') \]
      \item Parada: $\norminf{V_k - V_{k-1}} \le \varepsilon$ \quad
            {\small Custo por iteração: $O(\abs{\gS}^2)$}
    \end{itemize}
  \end{algbox}
  \vspace{-0.9em}
  \begin{columns}[T]
    \begin{column}{0.5\textwidth}
      \begin{ideiabox}
        $V_k(s)$ é o valor \emph{exato} de $s$ quando restam exatamente $k$ passos.
      \end{ideiabox}
    \end{column}
    \begin{column}{0.48\textwidth}
      \begin{alertabox}
        Analítico $\Theta(N^3)$ uma vez \emph{versus} iterativo $O(N^2)$ por iteração.
        Para $N$ grande o iterativo vence --- e é ele que sobrevive ao caso sem modelo.
      \end{alertabox}
    \end{column}
  \end{columns}
\end{frame}

%% ------------------------------------------------------------
\begin{frame}{Processo de decisão de Markov (MDP)}
  \begin{center}
    {\large MDP \;=\; processo de recompensa de Markov \;\alert{$+$ ações}}
  \end{center}
  \begin{defbox}{}
    Um MDP é a tupla $(\gS, \gA, P, \rec, \gamma)$:
    \begin{itemize}\setlength{\itemsep}{0.1ex}
      \item $\gS$: conjunto finito de estados de Markov; \quad $\gA$: conjunto finito de ações;
      \item $P$: modelo de transição, um por ação --- $\tran{s_{t+1}=s'}{s_t=s,\, a_t=a}$;
      \item $\rec$: recompensa --- $\rec(s_t=s, a_t=a) = \E[\, r_t \mid s_t=s,\, a_t=a \,]$;
      \item $\gamma \in [0,1]$: fator de desconto.
    \end{itemize}
  \end{defbox}
  \begin{alertabox}
    A recompensa aparece na literatura como $\rec(s)$, $\rec(s,a)$ ou $\rec(s,a,s')$.
    Usaremos \alert{$\rec(s,a)$}. As três têm o mesmo poder expressivo, mas trocar de
    convenção no meio de uma dedução é fonte clássica de erro.
  \end{alertabox}
\end{frame}

%% ------------------------------------------------------------
\begin{frame}{Exemplo: o MDP do \emph{Mars Rover}}
  \begin{center}
    \scalebox{0.84}{\marsroverdet}
  \end{center}
  \begin{columns}[T]
    \begin{column}{0.60\textwidth}
      \scalebox{0.68}{$
        P(s' \mid s, a_1) =
        \begin{bmatrix}
          1&0&0&0&0&0&0\\ 1&0&0&0&0&0&0\\ 0&1&0&0&0&0&0\\ 0&0&1&0&0&0&0\\
          0&0&0&1&0&0&0\\ 0&0&0&0&1&0&0\\ 0&0&0&0&0&1&0
        \end{bmatrix}
        \;
        P(s' \mid s, a_2) =
        \begin{bmatrix}
          0&1&0&0&0&0&0\\ 0&0&1&0&0&0&0\\ 0&0&0&1&0&0&0\\ 0&0&0&0&1&0&0\\
          0&0&0&0&0&1&0\\ 0&0&0&0&0&0&1\\ 0&0&0&0&0&0&1
        \end{bmatrix}
      $}
    \end{column}
    \begin{column}{0.38\textwidth}
      {\small
      \begin{itemize}\setlength{\itemsep}{0.25ex}
        \item $7$ estados: posição do rover.
        \item $2$ ações \alert{determinísticas}.
        \item Nas bordas, a ação que sairia do mapa mantém o rover parado.
        \item Recompensa: $+1$ em $s_1$, $+10$ em $s_7$, $0$ nos demais.
      \end{itemize}}
    \end{column}
  \end{columns}
\end{frame}

%% ------------------------------------------------------------
\begin{frame}{Políticas}
  \begin{defbox}{}
    Uma política especifica que ação tomar em cada estado. Na forma geral, é uma
    distribuição condicional: $\;\pol(a \mid s) = \Prob(a_t = a \mid s_t = s)$.
  \end{defbox}
  \begin{columns}[T]
    \begin{column}{0.49\textwidth}
      {\small
      \begin{itemize}\setlength{\itemsep}{0.4ex}
        \item \textbf{Determinística:} $\pol(s) = a$ --- massa concentrada em uma ação.
        \item \textbf{Estocástica:} massa distribuída entre ações.
        \item \textbf{Estacionária:} não depende de $t$.
      \end{itemize}}
    \end{column}
    \begin{column}{0.49\textwidth}
      \begin{ideiabox}
        A forma estocástica permite exploração, é diferenciável nos parâmetros
        (gradiente de política) e é a única ótima em jogos e em MDPs parcialmente
        observáveis.
      \end{ideiabox}
    \end{column}
  \end{columns}
  \begin{alertabox}
    Uma política \emph{não} é um plano. Plano é sequência fixa de ações; política é
    \alert{função reativa} do estado --- e por isso robusta à estocasticidade.
  \end{alertabox}
\end{frame}

%% ------------------------------------------------------------
\begin{frame}{MDP $+$ política $=$ MRP}
  Fixar uma política elimina a decisão: sobram aleatoriedade e recompensa.
  \[
    \rec^{\pol}(s) = \sum_{a \in \gA} \pol(a \mid s)\, \rec(s,a),
    \qquad
    P^{\pol}(s' \mid s) = \sum_{a \in \gA} \pol(a \mid s)\, \tran{s'}{s,a}.
  \]
  \begin{ideiabox}
    Esta é a observação mais rentável da aula. \alert{Tudo} que sabemos calcular para
    MRPs --- forma matricial, solução $(I - \gamma P^{\pol})^{-1} \rec^{\pol}$, algoritmo
    iterativo $O(\abs{\gS}^2)$ --- se aplica de imediato à \emph{avaliação} de uma
    política. O problema de predição não exige teoria nova.
  \end{ideiabox}
  \begin{alertabox}
    O que não vem de graça é o \alert{controle}: escolher $\pol$. É aí que entra o
    operador $\max$, e com ele a não linearidade.
  \end{alertabox}
\end{frame}

%% ============================================================
\section{Avaliação de política (predição)}
%% ============================================================

\begin{frame}{Avaliação iterativa de política em um MDP}
  \begin{algbox}{avaliação de política}
    \begin{itemize}\setlength{\itemsep}{0.2ex}
      \item Inicialize $V_0^{\pol}(s) = 0$ para todo $s$
      \item Para $k = 1, 2, \dots$ até convergir, e para todo $s \in \gS$:
        \[
          V_k^{\pol}(s) \leftarrow \sum_{a \in \gA} \pol(a \mid s)
          \left[\, \rec(s,a) + \gamma \sum_{s' \in \gS} \tran{s'}{s,a}\, V_{k-1}^{\pol}(s') \,\right]
        \]
    \end{itemize}
  \end{algbox}
  \vspace{-0.8em}
  {\small Este é o \alert{backup de Bellman para uma política específica}. Se $\pol$ for
  determinística, a soma externa desaparece:}
  \vspace{-0.4em}
  \[
    \textstyle V_k^{\pol}(s) \leftarrow \rec(s, \pol(s)) + \gamma \sum_{s' \in \gS} \tran{s'}{s,\pol(s)}\, V_{k-1}^{\pol}(s')
  \]
  \vspace{-0.9em}
  \begin{alertabox}
    Não há $\max$ aqui: a política já foi fixada. O operador é \alert{linear} em $V$ ---
    e é isso que permite resolver o sistema de forma fechada.
  \end{alertabox}
\end{frame}

%% ------------------------------------------------------------
\begin{frame}{Exercício L2E1 --- uma iteração de avaliação}
  \begin{center}\scalebox{0.78}{\marsroverdet}\end{center}
  \vspace{-0.3em}
  \begin{exbox}{--- \emph{Mars Rover}}
    Suponha agora que $a_1$ seja \alert{estocástica}:
    $p(s_6 \mid s_6, a_1) = 0{,}5$ e $p(s_7 \mid s_6, a_1) = 0{,}5$ (e assim por diante).

    \smallskip
    Recompensa: para toda ação, $+1$ em $s_1$, $+10$ em $s_7$, $0$ nos demais.
    Seja $\pol(s) = a_1$ para todo $s$, com
    $V_k = [\,1\;\,0\;\,0\;\,0\;\,0\;\,0\;\,10\,]$, $k=1$ e $\gamma = 0{,}5$.

    \smallskip
    \textbf{Calcule $V_{k+1}(s_6)$.}
  \end{exbox}
\end{frame}

%% ------------------------------------------------------------
\begin{frame}{Exercício L2E1 --- resolução}
  Aplicamos o backup para política determinística, em um único estado:
  \[
    V_{k+1}(s_6) = \rec(s_6, \pol(s_6))
      + \gamma \sum_{s' \in \gS} p(s' \mid s_6, a_1)\, V_k(s')
  \]
  Somente dois termos são não nulos, pois $p(\cdot \mid s_6, a_1)$ tem massa apenas em
  $s_6$ e $s_7$:
  \[
    \begin{aligned}
      V_{k+1}(s_6) &= \underbrace{0}_{\rec(s_6,a_1)}
        + \underbrace{0{,}5}_{\gamma}\cdot\bigl(\,0{,}5 \cdot \underbrace{0}_{V_k(s_6)}
        + 0{,}5 \cdot \underbrace{10}_{V_k(s_7)}\,\bigr) \\[0.2em]
      &= \mdestaque{rlVerde}{2{,}5}
    \end{aligned}
  \]
  \begin{ideiabox}
    Valor \emph{difunde} para trás pelas transições: na próxima varredura, $s_5$ também
    deixará de valer zero. Cada rodada propaga informação um passo --- por isso o número
    de iterações depende do diâmetro do MDP.
  \end{ideiabox}
\end{frame}

%% ------------------------------------------------------------
\begin{frame}{Teste sua compreensão --- o espaço de políticas}
  \begin{center}\scalebox{0.74}{\marsroverdet}\end{center}
  \vspace{-0.3em}
  \begin{quizbox}{}
    Em breve não bastará avaliar \emph{uma} política: queremos encontrar a \emph{melhor}.
    Vale a pena, então, entender o espaço de busca.
    \begin{enumerate}\setlength{\itemsep}{0.2ex}
      \item No \emph{Mars Rover} --- $7$ estados discretos, $2$ ações --- quantas
            políticas \alert{determinísticas} existem?
      \item A política ótima de um MDP (a de maior valor) é \alert{única}?
    \end{enumerate}
  \end{quizbox}
\end{frame}

%% ------------------------------------------------------------
\begin{frame}{Teste sua compreensão --- respostas}
  \begin{respbox}
    \textbf{1.} $2^7 = 128$. Em geral, há $\abs{\gA}^{\abs{\gS}}$ políticas
    determinísticas: uma escolha independente de ação para cada estado.

    \smallskip
    \textbf{2.} \textbf{Não.} Pode haver duas (ou infinitas) políticas com a
    \emph{mesma} função valor máxima --- basta que em algum estado dois valores $Q$
    empatem no máximo. O que é único é a \alert{função valor ótima} $\Vstar$, não o
    argumento que a atinge.
  \end{respbox}
  \begin{ideiabox}
    Distinção que vale guardar: $\Vstar$ é única, $\polstar$ não é. Isso reaparece o curso
    inteiro --- por exemplo, ao comparar duas execuções que convergiram para políticas
    diferentes com desempenho idêntico.
  \end{ideiabox}
\end{frame}

%% ------------------------------------------------------------
\begin{frame}{Três formas de avaliar uma política}
  \begin{center}
  {\small
  \begin{tabular}{@{}p{3.0cm}p{3.4cm}p{3.0cm}p{3.4cm}@{}}
    \toprule
    \textbf{Método} & \textbf{Custo} & \textbf{Exige Markov?} & \textbf{Exige modelo?} \\
    \midrule
    Analítico\newline $(I-\gamma P^{\pol})^{-1}\rec^{\pol}$
      & $\Theta(N^3)$, uma vez & sim & sim \\[0.4em]
    Iterativo (DP)
      & $O(N^2)$ por varredura & sim & sim \\[0.4em]
    Simulação\newline (Monte Carlo)
      & $O(\text{episódios} \times H)$ & \alert{não} & não (só amostras) \\
    \bottomrule
  \end{tabular}}
  \end{center}
  \vspace{0.3em}
  \begin{ideiabox}
    A terceira linha é a ponte para a próxima aula: gere muitos episódios e tire a média
    dos retornos. Desigualdades de concentração (Hoeffding) limitam a rapidez com que a
    média se aproxima da esperança. Não é preciso supor estrutura markoviana --- e,
    crucialmente, não é preciso conhecer $P$ nem $\rec$.
  \end{ideiabox}
\end{frame}

%% ============================================================
\section{Controle: iteração de política}
%% ============================================================

\begin{frame}{O problema de controle}
  \vspace{-0.9em}
  \[
    \polstar(s) = \argmax_{\pol} \Vpi(s)
  \]
  \vspace{-1.2em}
  \begin{columns}[T]
    \begin{column}{0.52\textwidth}
      {\small Para um MDP de horizonte infinito com $\gamma<1$, a política ótima é:
      \begin{itemize}\setlength{\itemsep}{0.2ex}
        \item \textbf{determinística} --- aleatorizar não ajuda;
        \item \textbf{estacionária} --- não depende de $t$;
        \item \textbf{não necessariamente única} --- pode haver empates.
      \end{itemize}
      \smallskip
      Já a \alert{função valor ótima} $\Vstar$ é única.}
    \end{column}
    \begin{column}{0.46\textwidth}
      \begin{alertabox}
        Leia o $\argmax$ com cuidado: ele afirma que existe uma política que maximiza
        o valor em \alert{todos os estados ao mesmo tempo}. Em otimização multiobjetivo
        isso em geral não existe; aqui existe, pela estrutura de Bellman.
      \end{alertabox}
    \end{column}
  \end{columns}
  \vspace{-0.2em}
  \begin{ideiabox}
    Por que estacionária? Com horizonte infinito, ``o que resta do problema'' é o mesmo em
    $t=3$ e em $t=3000$. Em horizonte finito isso falha --- voltaremos a esse ponto.
  \end{ideiabox}
\end{frame}

%% ------------------------------------------------------------
\begin{frame}{Buscar por enumeração é inviável}
  \begin{columns}[T]
    \begin{column}{0.5\textwidth}
      Uma opção seria enumerar todas as políticas determinísticas e avaliar cada uma.
      \begin{itemize}\setlength{\itemsep}{0.3ex}
        \item Número de políticas: $\abs{\gA}^{\abs{\gS}}$.
        \item Custo de avaliar cada uma: $\Theta(N^3)$ ou várias varreduras $O(N^2)$.
      \end{itemize}
      \vspace{0.4em}
      A \alert{iteração de política} faz muito melhor: ela usa o valor da política atual
      para construir diretamente uma política melhor, em vez de testar candidatas às cegas.
    \end{column}
    \begin{column}{0.48\textwidth}
      \begin{center}\small
      \begin{tabular}{@{}lll@{}}
        \toprule
        $\abs{\gS}$ & $\abs{\gA}$ & políticas \\
        \midrule
        $7$   & $2$ & $128$ \\
        $20$  & $4$ & $\approx 10^{12}$ \\
        $100$ & $4$ & $\approx 10^{60}$ \\
        \bottomrule
      \end{tabular}
      \end{center}
      \vspace{0.5em}
      \begin{alertabox}
        $10^{60}$ excede em muito o número de átomos da Via Láctea. Enumeração não é
        ``lenta'': é impossível.
      \end{alertabox}
    \end{column}
  \end{columns}
\end{frame}

%% ------------------------------------------------------------
\begin{frame}{Nova definição: valor de estado-ação $Q$}
  \begin{columns}[T]
    \begin{column}{0.62\textwidth}
      \begin{defbox}[fontupper=\small]{valor de estado-ação de uma política}
        \vspace{-1.1em}
        \[
          Q^{\pol}(s,a) = \rec(s,a) + \gamma \sum_{s' \in \gS} \tran{s'}{s,a}\, V^{\pol}(s')
        \]
        \vspace{-1.3em}

        ``Tome a ação $a$ \alert{agora} e siga $\pol$ \alert{daí em diante}.''
      \end{defbox}
      \vspace{-1.0em}
      \[
        \textstyle \Vpi(s) = \sum_{a \in \gA} \pol(a \mid s)\, Q^{\pol}(s,a)
      \]
      \vspace{-1.4em}
      \begin{ideiabox}
        $V$ diz ``quão bom é este estado?''; $Q$ diz ``quão boa é esta \emph{ação} aqui?''
        --- que é a pergunta necessária para \emph{agir}. Daí o domínio de $Q$ nos
        algoritmos sem modelo.
      \end{ideiabox}
    \end{column}
    \begin{column}{0.36\textwidth}
      \begin{center}\scalebox{0.78}{\backupotimo}\end{center}
      \vspace{-0.3em}
      {\scriptsize\color{rlCinza}\centering
      \emph{Backup}: um passo de antecipação a partir de $s$.\par}
    \end{column}
  \end{columns}
\end{frame}

%% ------------------------------------------------------------
\begin{frame}{Passo de melhoria de política}
  \begin{algbox}{melhoria de política (\emph{policy improvement})}
    \begin{itemize}\setlength{\itemsep}{0.3ex}
      \item Calcule o valor de estado-ação da política atual $\pol_i$: para todo
            $s \in \gS$, $a \in \gA$,
        \[
          Q^{\pol_i}(s,a) = \rec(s,a) + \gamma \sum_{s' \in \gS} \tran{s'}{s,a}\, V^{\pol_i}(s')
        \]
      \item Construa a nova política, gulosa em relação a $Q^{\pol_i}$:
        \[ \textstyle \pol_{i+1}(s) = \argmax_{a} Q^{\pol_i}(s,a) \quad \forall s \in \gS \]
    \end{itemize}
  \end{algbox}
  \begin{ideiabox}
    Nada garante \emph{a priori} que $\pol_{i+1}$ seja melhor: ela foi construída com
    $V^{\pol_i}$, o valor da política \alert{antiga}. Provar que a troca compensa é o
    teorema da melhoria de política.
  \end{ideiabox}
\end{frame}

%% ------------------------------------------------------------
\begin{frame}{Algoritmo: iteração de política (PI)}
  \begin{columns}[T]
    \begin{column}{0.61\textwidth}
      \begin{algbox}{Iteração de Política}
        \begin{itemize}\setlength{\itemsep}{0.25ex}
          \item $i \leftarrow 0$; inicialize $\pol_0$ aleatoriamente
          \item Enquanto $i = 0$ ou $\norm{\pol_i - \pol_{i-1}}_1 > 0$:
            \begin{itemize}\setlength{\itemsep}{0.15ex}
              \item $V^{\pol_i} \leftarrow$ \emph{avaliação} de $\pol_i$
              \item $\pol_{i+1} \leftarrow$ \emph{melhoria} sobre $V^{\pol_i}$
              \item $i \leftarrow i + 1$
            \end{itemize}
        \end{itemize}
      \end{algbox}
      {\small A norma $L_1$ só verifica se a política mudou em \emph{algum} estado.}
    \end{column}
    \begin{column}{0.38\textwidth}
      \begin{center}\scalebox{0.66}{\ciclogpi}\end{center}
      \vspace{-0.2em}
      {\scriptsize\color{rlCinza}\centering
      Os dois processos se perseguem até o ponto fixo $(\Vstar, \polstar)$.\par}
    \end{column}
  \end{columns}
  \vspace{-0.3em}
  \begin{ideiabox}
    Alternar avaliação e melhoria é a \alert{iteração de política generalizada},
    esqueleto de quase todo algoritmo do curso.
  \end{ideiabox}
\end{frame}

%% ------------------------------------------------------------
\begin{frame}{Por que a melhoria funciona: a intuição}
  Para qualquer estado $s$, escolher a melhor ação é pelo menos tão bom quanto seguir
  a política antiga:
  \[
    \max_{a} Q^{\pol_i}(s,a) \;\ge\; Q^{\pol_i}\bigl(s, \pol_i(s)\bigr)
    \;=\; \rec(s,\pol_i(s)) + \gamma \sum_{s'} \tran{s'}{s,\pol_i(s)} V^{\pol_i}(s')
    \;=\; V^{\pol_i}(s)
  \]
  \begin{itemize}\setlength{\itemsep}{0.3ex}
    \item Suponha que tomemos $\pol_{i+1}(s)$ \alert{por um passo} e depois sigamos
          $\pol_i$ para sempre: a soma esperada de recompensas é ao menos tão boa
          quanto sempre seguir $\pol_i$.
    \item \alert{Mas} a política proposta é seguir $\pol_{i+1}$ \emph{sempre}, não só
          por um passo.
  \end{itemize}
  \begin{alertabox}
    O argumento de um passo não fecha a prova. É preciso aplicá-lo repetidamente ---
    um telescópio de desigualdades --- para concluir que $V^{\pol_{i+1}} \ge V^{\pol_i}$.
  \end{alertabox}
\end{frame}

%% ------------------------------------------------------------
\begin{frame}{Melhoria monotônica de política}
  \begin{defbox}[fontupper=\small]{ordem parcial entre políticas}
    $V^{\pol_1} \ge V^{\pol_2}$ significa $V^{\pol_1}(s) \ge V^{\pol_2}(s)$
    para \alert{todo} $s \in \gS$.
  \end{defbox}
  \begin{teobox}{Proposição --- melhoria monotônica}
    Seja $\pol_{i+1}$ a política obtida pelo passo de melhoria aplicado a $\pol_i$. Então
    $V^{\pol_{i+1}} \ge V^{\pol_i}$, com desigualdade \alert{estrita} em ao menos um
    estado se $\pol_i$ for subótima.
  \end{teobox}
  \begin{ideiabox}
    Eis a resposta à pergunta que abriu a aula: PI \emph{de fato} melhora monotonicamente
    a cada rodada. Nem todo algoritmo tem essa garantia --- iteração de valor não a tem
    para as políticas intermediárias.
  \end{ideiabox}
\end{frame}

%% ------------------------------------------------------------
\begin{frame}{Prova: melhoria monotônica}
  \begin{center}
  \resizebox{0.99\textwidth}{!}{$\displaystyle
  \begin{aligned}
    V^{\pol_i}(s)
      &\le \max_{a} Q^{\pol_i}(s,a)
       = \max_{a} \Bigl[\, \rec(s,a) + \gamma \sum_{s'} \tran{s'}{s,a}\, V^{\pol_i}(s') \,\Bigr]
        && \text{\color{rlCinza} (i) o máximo domina} \\[0.2em]
      &= \rec(s, \pol_{i+1}(s)) + \gamma \sum_{s'} \tran{s'}{s,\pol_{i+1}(s)}\, V^{\pol_i}(s')
        && \text{\color{rlCinza} (ii) definição de $\pol_{i+1}$} \\[0.2em]
      &\le \rec(s, \pol_{i+1}(s)) + \gamma \sum_{s'} \tran{s'}{s,\pol_{i+1}(s)}
        \Bigl[\, \max_{a'} Q^{\pol_i}(s',a') \,\Bigr]
        && \text{\color{rlCinza} (iii) aplica (i) em $s'$} \\[0.2em]
      &= \rec(s, \pol_{i+1}(s)) + \gamma \sum_{s'} \tran{s'}{s,\pol_{i+1}(s)}
        \Bigl[\, \rec(s',\pol_{i+1}(s')) + \gamma \sum_{s''} \tran{s''}{s',\pol_{i+1}(s')}\, V^{\pol_i}(s'') \,\Bigr]
        && \text{\color{rlCinza} (iv) aplica (ii) em $s'$} \\[0.2em]
      &\;\;\vdots && \text{\color{rlCinza} repetir indefinidamente} \\[0.2em]
      &= V^{\pol_{i+1}}(s)
  \end{aligned}$}
  \end{center}
  \vspace{-0.4em}
  \begin{alertabox}
    A cada linha, mais um passo passa a ser executado sob $\pol_{i+1}$ e a cauda continua
    sob $\pol_i$. No limite tudo é $\pol_{i+1}$, e o resíduo some porque $\gamma^k \to 0$.
  \end{alertabox}
\end{frame}

%% ------------------------------------------------------------
\begin{frame}{Teste sua compreensão --- iteração de política}
  \begin{algbox}{Iteração de Política (espaços finitos)}
    \begin{itemize}\setlength{\itemsep}{0.2ex}
      \item $i \leftarrow 0$; inicialize $\pol_0$ aleatoriamente
      \item Enquanto $i=0$ ou $\norm{\pol_i - \pol_{i-1}}_1 > 0$: avaliar $\pol_i$,
            melhorar para $\pol_{i+1}$, $i \leftarrow i+1$
    \end{itemize}
  \end{algbox}
  \begin{quizbox}{}
    \begin{enumerate}\setlength{\itemsep}{0.3ex}
      \item Se a política \alert{não muda} em uma rodada, ela pode voltar a mudar depois?
      \item Existe um \alert{número máximo} de iterações de iteração de política?
    \end{enumerate}
  \end{quizbox}
\end{frame}

%% ------------------------------------------------------------
\begin{frame}{Teste sua compreensão --- respostas}
  \begin{respbox}
    \textbf{1. Não.} Se $\pol_{i+1}(s) = \pol_i(s)$ para todo $s$, então
    $Q^{\pol_{i+1}} = Q^{\pol_i}$ (mesma política, mesmo valor). Logo
    \[
      \pol_{i+2}(s) = \argmax_a Q^{\pol_{i+1}}(s,a) = \argmax_a Q^{\pol_i}(s,a) = \pol_{i+1}(s).
    \]
    O algoritmo atingiu um ponto fixo, e esse ponto fixo é ótimo.
  \end{respbox}
  \begin{respbox}
    \textbf{2. Sim: $\abs{\gA}^{\abs{\gS}}$.} Pela monotonia, cada política só aparece em
    \emph{uma} rodada (a menos que já seja ótima) --- e esse é o número de políticas
    determinísticas.
  \end{respbox}
  \begin{ideiabox}
    Na prática PI converge em poucas dezenas de iterações; Ye (2011) mostrou que, para
    $\gamma$ fixo, PI é \alert{fortemente polinomial}.
  \end{ideiabox}
\end{frame}

%% ============================================================
\section{Controle: iteração de valor}
%% ============================================================

\begin{frame}{Uma segunda estratégia de controle}
  \begin{columns}[T]
    \begin{column}{0.49\textwidth}
      {\color{rlAzulClaro}\small\bfseries ITERAÇÃO DE POLÍTICA}\\[0.3em]
      {\small Mantém uma \alert{política} e calcula seu valor de horizonte
      \emph{infinito}; usa esse valor para propor uma política melhor.}
    \end{column}
    \begin{column}{0.49\textwidth}
      {\color{rlLaranja}\small\bfseries ITERAÇÃO DE VALOR}\\[0.3em]
      {\small Mantém o \alert{valor ótimo} de começar em cada estado supondo que restam
      $k$ passos de episódio --- e itera considerando episódios cada vez mais longos.}
    \end{column}
  \end{columns}
  \vspace{0.6em}
  \begin{ideiabox}
    São dois eixos de aproximação ortogonais: PI é exata no \emph{valor} e aproxima a
    \emph{política}; VI é exata na \emph{otimalidade local} e aproxima o \emph{horizonte}.
    Ambas convergem para o mesmo par $(\Vstar, \polstar)$.
  \end{ideiabox}
  \begin{alertabox}
    Em VI, os valores intermediários $V_k$ em geral \alert{não} correspondem ao valor de
    nenhuma política executável --- são valores ótimos de um problema truncado.
  \end{alertabox}
\end{frame}

%% ------------------------------------------------------------
\begin{frame}{Equação de Bellman e operador de backup}
  \begin{defbox}[fontupper=\small]{equação de Bellman (para uma política)}
    O valor de uma política deve satisfazer
    \(\; \Vpi(s) = \rec^{\pol}(s) + \gamma \sum_{s' \in \gS} P^{\pol}(s' \mid s)\, \Vpi(s'). \)
  \end{defbox}
  \begin{defbox}[fontupper=\small]{operador de backup de Bellman $\Bop$}
    Devolve uma nova função valor sobre \emph{todos} os estados, melhorando o valor
    quando há melhora possível:
    \vspace{-0.5em}
    \[
      \Bop V(s) = \max_{a} \Bigl[\, \rec(s,a) + \gamma \sum_{s' \in \gS} \tran{s'}{s,a}\, V(s') \,\Bigr]
    \]
    \vspace{-1.4em}
  \end{defbox}
  \begin{ideiabox}
    A troca da \emph{soma ponderada por $\pol$} pelo \emph{$\max$ sobre ações} é o que
    distingue predição de controle --- e o que destrói a linearidade: com $\Bop$, a
    solução fechada por inversão de matriz deixa de existir.
  \end{ideiabox}
\end{frame}

%% ------------------------------------------------------------
\begin{frame}{Algoritmo: iteração de valor (VI)}
  \begin{algbox}{Iteração de Valor}
    \begin{itemize}\setlength{\itemsep}{0.25ex}
      \item $k \leftarrow 1$; inicialize $V_0(s) = 0$ para todo $s$
      \item Repita até convergir (por ex., $\norminf{V_{k+1} - V_k} \le \epsilon$):
            para cada $s \in \gS$,
        \[
          V_{k+1}(s) = \max_{a} \Bigl[\, \rec(s,a) + \gamma \sum_{s' \in \gS} \tran{s'}{s,a}\, V_k(s') \,\Bigr]
          \qquad \text{ou seja,} \quad V_{k+1} = \Bop V_k
        \]
    \end{itemize}
  \end{algbox}
  \vspace{-0.4em}
  {\small Ao final, extraia a política \alert{gulosa} em relação ao último valor:}
  \vspace{-0.3em}
  \[
    \textstyle \pol(s) = \argmax_{a} \bigl[\, \rec(s,a) + \gamma \sum_{s'} \tran{s'}{s,a}\, V_{k+1}(s') \,\bigr]
  \]
  \vspace{-1.0em}
  \begin{ideiabox}
    Leitura em horizonte: $V_k$ é o valor ótimo se restarem $k$ passos; a política gulosa
    extraída é a ação ótima quando se pode agir por $k+1$ passos.
  \end{ideiabox}
\end{frame}

%% ------------------------------------------------------------
\begin{frame}{Iteração de política na linguagem de operadores}
  O operador de backup \alert{para uma política fixa} $\pol$ é
  \[
    \Bpol V(s) = \rec^{\pol}(s) + \gamma \sum_{s' \in \gS} P^{\pol}(s' \mid s)\, V(s')
  \]
  \begin{columns}[T]
    \begin{column}{0.49\textwidth}
      \textbf{Avaliação de política} = calcular o \alert{ponto fixo} de $\Bpol$:
      aplicar o operador até $V$ parar de mudar,
      \[
        \Vpi = \Bpol \Bpol \cdots \Bpol V.
      \]
    \end{column}
    \begin{column}{0.49\textwidth}
      \textbf{Melhoria de política} = um backup com $\max$, seguido de $\argmax$:
      \vspace{-0.3em}
      \begin{center}
      \resizebox{0.98\linewidth}{!}{$
        \pol_{k+1}(s) = \argmax_a \Bigl[ \rec(s,a) + \gamma \sum_{s'} \tran{s'}{s,a} V^{\pol_k}(s') \Bigr]
      $}
      \end{center}
    \end{column}
  \end{columns}
  \vspace{0.4em}
  \begin{ideiabox}
    Na mesma notação: VI itera $\Bop$ (com $\max$) desde o início; PI itera $\Bpol$
    (linear) até o ponto fixo e só então usa o $\max$ uma vez. Essa lente unificada
    prepara o terreno para variantes híbridas.
  \end{ideiabox}
\end{frame}

%% ------------------------------------------------------------
\begin{frame}{Operador de contração}
  \begin{defbox}[fontupper=\small]{operador de contração}
    Seja $O$ um operador e $\norm{\cdot}$ uma norma. Se
    $\norm{OV - OV'} \le \norm{V - V'}$ para todos $V, V'$, então $O$ é um operador de
    \emph{contração} (estrita, se vale com fator $\gamma<1$).
  \end{defbox}
  \begin{ideiabox}
    Imagem mental: um mapa que \alert{aproxima os pontos}. Aplicado repetidamente,
    todos colapsam em um único ponto fixo --- não importa de onde se parta.
  \end{ideiabox}
  \begin{teobox}{VI converge?}
    Sim, se $\gamma < 1$ --- ou se um estado terminal é atingido com probabilidade $1$.
    O backup de Bellman é uma contração para $\gamma<1$: aplicado a duas funções valor
    distintas, a distância entre elas \alert{encolhe} por um fator $\gamma$.
  \end{teobox}
\end{frame}

%% ------------------------------------------------------------
\begin{frame}{Prova: o backup de Bellman é uma contração ($\gamma<1$)}
  Norma do supremo: $\norminf{V - V'} = \max_s \abs{V(s) - V'(s)}$. Para cada $s$:
  \begin{center}
  \resizebox{0.97\textwidth}{!}{$\displaystyle
  \begin{aligned}
    \abs{\Bop V_k(s) - \Bop V_j(s)}
      &= \Bigl|\, \max_a \Bigl[ \rec(s,a) + \gamma \sum_{s'} \tran{s'}{s,a} V_k(s') \Bigr]
         - \max_{a'} \Bigl[ \rec(s,a') + \gamma \sum_{s'} \tran{s'}{s,a'} V_j(s') \Bigr] \,\Bigr| \\[0.15em]
      &\le \max_a \Bigl|\, \rec(s,a) + \gamma \sum_{s'} \tran{s'}{s,a} V_k(s')
         - \rec(s,a) - \gamma \sum_{s'} \tran{s'}{s,a} V_j(s') \,\Bigr|
         && \text{\color{rlCinza} $\abs{\max f - \max g} \le \max \abs{f-g}$} \\[0.15em]
      &= \max_a \Bigl|\, \gamma \sum_{s'} \tran{s'}{s,a} \bigl( V_k(s') - V_j(s') \bigr) \,\Bigr|
         \le \max_a \, \gamma \sum_{s'} \tran{s'}{s,a}\, \norminf{V_k - V_j}
         && \text{\color{rlCinza} cota termo a termo} \\[0.15em]
      &= \gamma\, \norminf{V_k - V_j} \sum_{s'} \tran{s'}{s,a} = \gamma\, \norminf{V_k - V_j}
         && \text{\color{rlCinza} probabilidades somam $1$}
  \end{aligned}$}
  \end{center}
  Tomando o máximo sobre $s$:\; $\norminf{\Bop V_k - \Bop V_j} \le \gamma \norminf{V_k - V_j}$. \qed
  \vspace{0.1em}
  \begin{alertabox}
    Mesmo que todas as desigualdades fossem igualdades, ainda seria contração se
    $\gamma<1$. E a recompensa $\rec$ \alert{cancela} --- a contração não depende dela.
  \end{alertabox}
\end{frame}

%% ------------------------------------------------------------
\begin{frame}{Da contração à convergência}
  \begin{teobox}{ponto fixo de Banach aplicado a VI}
    Se $\gamma<1$, então: (i) $\Bop$ tem um \alert{único} ponto fixo, que é $\Vstar$;
    (ii) VI converge a partir de \alert{qualquer} inicialização $V_0$; e (iii) a
    convergência é geométrica:
    \(\; \norminf{V_k - \Vstar} \le \gamma^k\, \norminf{V_0 - \Vstar}. \)
  \end{teobox}
  \begin{columns}[T]
    \begin{column}{0.52\textwidth}
      {\small Consequências práticas:
      \begin{itemize}\setlength{\itemsep}{0.3ex}
        \item A inicialização não afeta o destino --- só o tempo de viagem (e uma boa
              inicialização \emph{acelera}).
        \item $\gamma$ perto de $1$ torna a convergência lenta: o número de iterações
              cresce como $\dfrac{1}{1-\gamma}$.
      \end{itemize}}
    \end{column}
    \begin{column}{0.46\textwidth}
      \begin{ideiabox}
        Parou com $\norminf{V_{k+1}-V_k} \le \epsilon$? Então
        $\norminf{V_{k+1} - \Vstar} \le \dfrac{\gamma\,\epsilon}{1-\gamma}$, e a política
        gulosa extraída perde no máximo $\dfrac{2\gamma\,\epsilon}{1-\gamma}$ em valor.
      \end{ideiabox}
    \end{column}
  \end{columns}
\end{frame}

%% ------------------------------------------------------------
\begin{frame}{Prática recomendada (fora de aula)}
  \begin{exbox}{--- para consolidar}
    \begin{enumerate}\setlength{\itemsep}{0.5ex}
      \item Prove que a iteração de valor converge para uma solução \emph{única} em
            espaços discretos de estados e ações com $\gamma < 1$.
      \item A inicialização dos valores em VI afeta alguma coisa? O quê exatamente ---
            e o que ela \emph{não} afeta?
      \item O valor da política extraída de VI a cada rodada é garantidamente
            \alert{monotônico} (se executada no problema real de horizonte infinito),
            como na iteração de política?
    \end{enumerate}
  \end{exbox}
  \begin{ideiabox}
    A terceira pergunta é a mais sutil da aula: a \emph{função valor} $V_k$ converge
    monotonicamente em norma, mas a sequência de \emph{políticas gulosas} extraídas pode
    oscilar em qualidade antes de estabilizar. Contraste direto com PI.
  \end{ideiabox}
\end{frame}

%% ============================================================
\section{Horizonte finito e síntese}
%% ============================================================

\begin{frame}{Iteração de valor com horizonte finito $H$}
  Em horizonte finito, VI deixa de ser aproximação e vira \alert{o algoritmo exato}:
  $V_k$ e $\pol_k$ são o valor e a política ótimos \emph{quando restam $k$ decisões}.
  \begin{algbox}{VI para horizonte $H$}
    \begin{itemize}\setlength{\itemsep}{0.2ex}
      \item Inicialize $V_0(s) = 0$ para todo $s$
      \item Para $k = 1, \dots, H$, e para cada $s$:
        \begin{center}
        \resizebox{0.96\linewidth}{!}{$
          V_{k+1}(s) = \max_a \Bigl[ \rec(s,a) + \gamma \sum_{s'} \tran{s'}{s,a} V_k(s') \Bigr],
          \qquad
          \pol_{k+1}(s) = \argmax_a \Bigl[ \rec(s,a) + \gamma \sum_{s'} \tran{s'}{s,a} V_k(s') \Bigr]
        $}
        \end{center}
    \end{itemize}
  \end{algbox}
  \vspace{-0.2em}
  \begin{ideiabox}
    Guarde \emph{todas} as $\pol_1, \dots, \pol_H$: para agir no instante $t$, use
    $\pol_{H-t}$ --- a política indexada por \alert{quantos passos restam}.
  \end{ideiabox}
\end{frame}

%% ------------------------------------------------------------
\begin{frame}{Avaliação por simulação (ponte para a próxima aula)}
  Alternativa para estimar o valor de uma política em horizonte finito:
  \begin{itemize}\setlength{\itemsep}{0.3ex}
    \item gere um grande número de \alert{episódios} seguindo a política;
    \item calcule o retorno de cada episódio e tire a \alert{média};
    \item desigualdades de concentração limitam a rapidez com que a média se aproxima da
          esperança.
  \end{itemize}
  \begin{ideiabox}
    Nenhuma suposição markoviana é necessária --- e, mais importante, nenhum
    \alert{modelo}: basta poder \emph{interagir}. Este é o portão de entrada do RL
    propriamente dito.
  \end{ideiabox}
  \begin{alertabox}
    O preço: variância. Cada retorno é uma amostra ruidosa, e a precisão melhora só com
    $1/\sqrt{n}$ episódios. O contraste exato entre DP e Monte Carlo é o tema da Aula 3.
  \end{alertabox}
\end{frame}

%% ------------------------------------------------------------
\begin{frame}{Exemplo: retornos amostrais no \emph{Mars Rover} estocástico}
  \begin{center}\scalebox{0.58}{\marsroverestoc}\end{center}\vspace{-0.6em}
  \vspace{-0.5em}
  {\small Recompensa $+1$ em $s_1$, $+10$ em $s_7$, $0$ nos demais. Episódios de $H=4$
  passos a partir de $s_4$, com $\gamma = \tfrac{1}{2}$:}
  \begin{center}
  \renewcommand{\arraystretch}{0.95}%
  {\footnotesize
  \begin{tabular}{@{}llc@{}}
    \toprule
    \textbf{Episódio} & \textbf{Retorno} $G = r_1 + \tfrac12 r_2 + \tfrac14 r_3 + \tfrac18 r_4$ & \textbf{Valor} \\
    \midrule
    $s_4, s_5, s_6, s_7$ & $0 + \tfrac12 \cdot 0 + \tfrac14 \cdot 0 + \tfrac18 \cdot 10$ & $1{,}25$ \\
    $s_4, s_4, s_5, s_4$ & $0 + \tfrac12 \cdot 0 + \tfrac14 \cdot 0 + \tfrac18 \cdot 0$  & $0$ \\
    $s_4, s_3, s_2, s_1$ & $0 + \tfrac12 \cdot 0 + \tfrac14 \cdot 0 + \tfrac18 \cdot 1$  & $0{,}125$ \\
    \bottomrule
  \end{tabular}}
  \end{center}
  \vspace{-0.5em}
  \begin{ideiabox}[fontupper=\footnotesize]
    Três amostras, três respostas --- todas legítimas. $V(s_4)$ é a média dessa
    distribuição; estimá-lo bem exige muitos episódios.
  \end{ideiabox}
\end{frame}

%% ------------------------------------------------------------
\begin{frame}{A política ótima é estacionária em horizonte finito?}
  \begin{quizbox}{}
    Em tarefas de horizonte \alert{finito}, a política ótima é estacionária
    (independente do passo de tempo)?
  \end{quizbox}
  \pause
  \begin{respbox}
    \textbf{Em geral, não.} A ação ótima depende de \alert{quanto tempo resta}.
  \end{respbox}
  \begin{ideiabox}
    Exemplo no \emph{Mars Rover}: em $s_4$, com muitos passos restantes vale a pena
    caminhar até o $+10$ em $s_7$; com \alert{um único} passo restante, nenhuma das duas
    direções alcança recompensa --- e com dois ou três, o $+1$ de $s_1$ e o $+10$ de
    $s_7$ competem de forma diferente a cada instante. A política ótima precisa ser
    indexada pelo tempo: $\pol_t^*(s)$.
  \end{ideiabox}
\end{frame}

%% ------------------------------------------------------------
\begin{frame}{Iteração de valor \emph{versus} iteração de política}
  \vspace{-0.3em}
  \begin{center}
  {\small
  \begin{tabular}{@{}p{3.4cm}p{4.6cm}p{4.6cm}@{}}
    \toprule
     & \textbf{Iteração de valor} & \textbf{Iteração de política} \\
    \midrule
    O que mantém & valor ótimo p/ horizonte $k$ & política + seu valor $\infty$ \\
    Passo & $V_{k+1} = \Bop V_k$ & avaliar $\Bpol$; melhorar c/ $\max$ \\
    Custo por iteração & $O(\abs{\gS}^2 \abs{\gA})$ & avaliação $+$ $O(\abs{\gS}^2 \abs{\gA})$ \\
    N\textsuperscript{o} de iterações & muitas e baratas & poucas e caras \\
    Monotonia da política & não garantida & \alert{garantida} \\
    Parente em RL moderno & Q-learning & gradiente de política \\
    \bottomrule
  \end{tabular}}
  \end{center}
  \vspace{-0.1em}
  \begin{ideiabox}
    \alert{Iteração de política modificada}: no passo de avaliação, aplique $\Bpol$ só
    $m$ vezes. Com $m=1$ recupera-se VI; com $m=\infty$, PI. Quase tudo entre os dois
    extremos funciona --- e costuma ser mais rápido que ambos.
  \end{ideiabox}
\end{frame}

%% ------------------------------------------------------------
\begin{frame}{Vocabulário de RL: modelos, políticas, valores}
  \begin{defbox}{modelo}
    Representação matemática da dinâmica ($P$) e da recompensa ($\rec$).
  \end{defbox}
  \begin{defbox}{política}
    Função que mapeia estados em ações (ou em distribuições sobre ações).
  \end{defbox}
  \begin{defbox}{função valor}
    Recompensas futuras esperadas a partir de um estado (e/ou ação), quando se segue uma
    política específica.
  \end{defbox}
  \begin{ideiabox}
    Estes três objetos são os eixos que organizam \emph{todo} o campo: algoritmos
    baseados em modelo, em valor ou em política --- e as combinações entre eles.
  \end{ideiabox}
\end{frame}

%% ------------------------------------------------------------
\begin{frame}{O que você deve saber ao final desta aula}
  \begin{columns}[T]
    \begin{column}{0.52\textwidth}
      {\small
      \begin{itemize}\setlength{\itemsep}{0.5ex}
        \item \textbf{Definir}: MP, MRP, MDP, operador de Bellman, contração, modelo,
              valor $Q$, política.
        \item \textbf{Implementar}: iteração de valor e iteração de política.
        \item \textbf{Comparar}: prós e contras das abordagens de avaliação de política.
      \end{itemize}}
    \end{column}
    \begin{column}{0.46\textwidth}
      {\small
      \begin{itemize}\setlength{\itemsep}{0.5ex}
        \item \textbf{Provar}: propriedades de contração e melhoria monotônica.
        \item \textbf{Criticar}: limitações das abordagens e da hipótese de Markov.
        \item \textbf{Distinguir}: quais métodos de avaliação exigem a hipótese de
              Markov --- e quais não.
      \end{itemize}}
    \end{column}
  \end{columns}
  \vspace{0.4em}
  \begin{alertabox}
    Pergunta de despedida, para levar à Aula 3: dos três métodos de avaliação vistos
    (analítico, programação dinâmica, simulação), quais exigem a hipótese de Markov?
  \end{alertabox}
\end{frame}

%% ------------------------------------------------------------
\begin{frame}[plain, noframenumbering]
  \begin{tikzpicture}[remember picture, overlay]
    \fill[rlAzul] (current page.south west) rectangle (current page.north east);
    \node[anchor=west, text width=0.8\paperwidth]
      at ([xshift=1.1cm]current page.west) {%
        {\color{rlAmbar}\small\bfseries PRÓXIMA AULA}\\[0.6em]
        {\color{white}\LARGE\bfseries Avaliação de política sem modelo}\\[0.6em]
        {\color{white!80!rlAzul}Monte Carlo \textbullet\ diferenças temporais \textbullet\
        o primeiro contato com aprendizagem de verdade}\\[1.4em]
        {\color{rlLaranja}\rule{3em}{2.4pt}}\\[1.2em]
        {\color{white!65!rlAzul}\footnotesize Material adaptado de Emma Brunskill,
        CS234: Reinforcement Learning, Stanford University.}};
  \end{tikzpicture}
\end{frame}

\end{document}
